\documentclass[12pt,letterpaper]{article}

\usepackage[
    top=2cm,
    bottom=2cm,
    left=2cm,
    right=2cm,
    footskip=1cm
]{geometry}
\usepackage{iftex}
\ifPDFTeX
    \input{glyphtounicode}
    \pdfgentounicode=1
    \usepackage[T1]{fontenc}
    \usepackage[utf8]{inputenc}
\fi
\usepackage{charter}
\usepackage{microtype}
\usepackage{array}
\usepackage{tabularx}
\usepackage{titlesec}
\usepackage{needspace}
\usepackage[dvipsnames]{xcolor}
\definecolor{primaryColor}{RGB}{0,0,0}
\usepackage[
    pdftitle={Zeqi Wu's CV},
    pdfauthor={Zeqi Wu},
    colorlinks=true,
    linkcolor=primaryColor,
    urlcolor=primaryColor
]{hyperref}
\usepackage{bookmark}

\providecommand{\LyX}{\texorpdfstring{%
    L\kern-.1667em\lower.25em\hbox{Y}\kern-.125emX}{LyX}}

\pagestyle{empty}
\setcounter{secnumdepth}{0}
\setlength{\parindent}{0pt}
\setlength{\parskip}{0pt}
\setlength{\tabcolsep}{0pt}
\linespread{1.06}
\raggedright

\titleformat{\section}
    {\needspace{4\baselineskip}\large\bfseries}
    {}{0pt}{}[\vspace{1pt}\titlerule]
\titlespacing*{\section}{0pt}{10pt}{6pt}

% Arguments: link, title, coauthors in alphabetical order, publication details.
\newcommand{\paperentry}[4]{%
    \par\addvspace{7pt}%
    \needspace{4\baselineskip}%
    \begingroup
    \hangindent=1.25em
    \hangafter=1
    \href{#1}{``#2''} (with #3), #4.\par
    \endgroup
}

\begin{document}

\begin{center}
    {\fontsize{25pt}{30pt}\selectfont Zeqi Wu}\par
    \vspace{5pt}
    \href{mailto:wuzeqi@ruc.edu.cn}{wuzeqi@ruc.edu.cn}
    \quad\textbar\quad
    \href{https://www.zeqiwu.com}{www.zeqiwu.com}
\end{center}
\vspace{-5pt}

\section{Education}

\begin{tabularx}{\textwidth}{@{}X r@{}}
    \textbf{Renmin University of China}
        & Sept.\ 2023 -- June 2028 (expected) \\
    Ph.D.\ in Statistics & \\[5pt]
    \textbf{Xiamen University}
        & Sept.\ 2019 -- June 2023 \\
    B.S.\ in Statistics &
\end{tabularx}

\section{Research Interests}

Causal inference; semiparametric and nonparametric estimation and inference;
machine learning; policy evaluation and learning; network econometrics.

\section{Peer-reviewed Publications}

\paperentry
    {https://doi.org/10.1017/S026646662400015X}
    {Applications of Functional Dependence to Spatial Econometrics}
    {Wen~Jiang and Xingbai~Xu}
    {\textit{Econometric Theory}, 41(5), 1044--1079,~2025}

\section{Working Papers}

\paperentry
    {https://arxiv.org/abs/2606.01659}
    {Data-Automated Policy Learning for Nonlinear Welfare}
    {Chunrong~Ai and Zheng~Zhang}
    {Working~Paper,~2026}

\paperentry
    {https://arxiv.org/abs/2511.17928}
    {Limit Theorems for Network Data without Metric Structure}
    {Wen~Jiang, Yachen~Wang, and Xingbai~Xu}
    {Working~Paper,~2026}

\paperentry
    {https://arxiv.org/abs/2603.07055}
    {Integrating Heterogeneous Information in Randomized Experiments:
     A Unified Calibration Framework}
    {Wei~Ma and Zheng~Zhang}
    {Working~Paper,~2026}

\paperentry
    {https://arxiv.org/abs/2507.04044}
    {A New and Efficient Debiased Estimation of General Treatment Models
     by Balanced Neural Networks Weighting}
    {Wei~Huang, Meilin~Wang, and Zheng~Zhang}
    {Working~Paper,~2025}

\section{Honors \& Awards}

\begin{tabularx}{\textwidth}{@{}X r@{}}
    \textbf{National Scholarship}, Xiamen University & Dec.\ 2021
\end{tabularx}

\section{Skills}

\textbf{Languages:} Mandarin (native), English\par
\textbf{Programming:} Python, R, MATLAB\par
\textbf{Typesetting:} \LaTeX, \LyX

\end{document}
